% New contribution after checkpoint08; root reserved A.4.48.
\subsection{A genuine harmonic survivor at the ternary signed five-point level}

We use the discrete marked source
\[
 B_0=\langle s,r\mid s^2=r^3=1\rangle,
 \qquad F_0=\langle x,y\rangle,\qquad
 x=(sr)^{-2},\quad y=sxs,
 \tag{HTS.1}
\]
and its actual profinite completion.  Put
$Q_B=\operatorname{PSL}_2(\mathbb Z/12)$, with the
simultaneous central quotient by $\{I,-I\}$, and retain
the marked epimorphism $q_B:B_0\to Q_B$ of A.4.45.
Thus $|Q_B|=576$.  Write $K=\ker q_B$.
On $F_0$, this projective map has the same kernel as
the signed map $\theta_{12}$: its signed image has
order $96$, and its image modulo four is the Klein
four group omitting $-I$.

Let $\Gamma=\Gamma_{0,5}$, let
$q:\Gamma\to Q_{192}$ be the actual signed quotient
of A.4.45, and put
\[
 N=\ker q,\qquad \Gamma_3=\Gamma/(N'N^3).
 \tag{HTS.2}
\]
The notation here denotes ordinary finite quotients;
each map extends uniquely to the respective profinite
completion.  Use the eight literal insertions
$\iota_i:F_0\to\Gamma$ in (SVM.1), and put
\[
 J_0=\bigcap_{i=1}^{8}\ker(F_0\xrightarrow{\iota_i}
                                     \Gamma_3),
 \qquad E=B_0/J_0,\qquad P=F_0/J_0.
 \tag{HTS.3}
\]
The exact source calculation (SVM.2)--(SVM.6) proves
that $J_0$ is normal in $B_0$.  The following are
properties of this one marked source, including its
actual quotient maps:
\[
 \begin{gathered}
 |E|=1728\cdot3^{97},\qquad
 E\twoheadrightarrow Q_B,\qquad
 |\ker(E\to Q_B)|=3^{98},\\
 P^{\mathrm{ab}}=C_{18}\overline x\oplus
                   C_{18}\overline y,
 \qquad |\overline x|=|\overline y|=18.
 \end{gathered}                                      \tag{HTS.4}
\]
For the last line, the base order-$288$ pure image has
derived subgroup of order eight and abelianization
$C_6^2$.  In the retained $97$-dimensional source
module, the divided exponent map
$(\operatorname{exp}_x/6,\operatorname{exp}_y/6)$
modulo three factors through the exact common
projection.  The sixth powers of $x,y$ map to its
two standard basis vectors.  Hence the exponent
lattice of $J_0$ is exactly $18\mathbb Z^2$;
also $x^{18},y^{18}$ are killed in every retained
insertion.  This proves the last line of (HTS.4),
including the named generators.  The coefficientwise
factorization and these power calculations are saved
in \path{check_fivepoint_signed_orders.py}, its output,
and its numerical archive.

\begin{result}{A.4.48}{Proved}{A full parity-one identity record at the ternary five-point level}
There exists an actual full multiplier-one gentle
element $g=(1,f_g)$ satisfying
\[
 \Delta(g)=1,\qquad (f_g)_P=1,
 \qquad (A_g)_P=1,
 \qquad \iota_i(f_g)=1\text{ in }\Gamma_3
                  \quad(1\leq i\leq8).
 \tag{HTS.5}
\]
Its normalized harmonic frame has a specified
nontrivial projective level-three harmonic value.
In particular the full pentagon word is one in
$\Gamma_3$ for this genuine source element.
The construction uses the identity automorphism of
the full canonical central cover of $E$, with a
nontrivial pure harmonic frame and cubic frame one.
\end{result}

\begin{proof}
We first identify the full canonical central kernels
and their actual comparison.  The source
abelianization is $B_0^{\mathrm{ab}}=C_6$.
The marked map to $Q_B^{\mathrm{ab}}=C_6$ is an
isomorphism, so
\[
                  J_0\subset K\subset B_0'.
 \tag{HTS.6}
\]
For $T=E,Q_B$, let $K_T$ be its kernel in $B_0$ and
set
\[
 T^c=B_0/[B_0,K_T],\qquad
 M_T=K_T/[B_0,K_T].
 \tag{HTS.7}
\]
Here $[B_0,K_T]$ means the normal subgroup generated
by all the indicated commutators.  The integral
homology five-term sequence
\cite[6.8.3, p.~196]{CA_WeibelHA1994} and
$H_2(B_0,\mathbb Z)=H_2(C_2,\mathbb Z)\oplus
H_2(C_3,\mathbb Z)=0$ give natural isomorphisms
\[
                       M_T=H_2(T,\mathbb Z).
 \tag{HTS.8}
\]
Indeed (HTS.6) makes the following map
$M_T\to B_0^{\mathrm{ab}}$ zero, so the preceding
homology map is onto as well as one-to-one.
These groups are finite, because $T$ is finite.
Consequently the displayed ordinary central covers
are finite and are exactly their profinite versions
$\widehat B_0/\overline{[\widehat B_0,\overline K_T]}$.
Thus no completion comparison is suppressed.

The inclusion $J_0\subset K$ gives the literal
surjection of marked covers
\[
 \begin{array}{ccc}
 E^c=B_0/[B_0,J_0]&\longrightarrow&Q_B^c=B_0/[B_0,K]\\
 \big\downarrow&&\big\downarrow\\
 E&\longrightarrow&Q_B .
 \end{array}                                        \tag{HTS.9}
\]
Equivalently (HTS.9) is the commutative square whose
horizontal and vertical arrows all send the two
named generators to the equally named generators.
Under (HTS.8), its map $M_E\to M_{Q_B}$ is the
induced integral homology map, by naturality of the
five-term sequence.

Put $L=\ker(E\to Q_B)$.  The integral homology
Lyndon--Hochschild--Serre spectral sequence
\cite[Theorem~6.8.2, pp.~195--196]{CA_WeibelHA1994} for
$1\to L\to E\to Q_B\to1$, localized at two, has
only its row of degree zero in $L$: every positive
$H_j(L,\mathbb Z)$ is annihilated by the odd integer
$|L|=3^{98}$
\cite[Theorem~6.5.8 and Corollary~6.5.10, p.~180]{CA_WeibelHA1994}.
Hence its edge map is the actual
isomorphism
\[
 H_2(E,\mathbb Z)\otimes\mathbb Z_{(2)}
     \xrightarrow{\ \sim\ }
 H_2(Q_B,\mathbb Z)\otimes\mathbb Z_{(2)}.
 \tag{HTS.10}
\]
For finite abelian groups localization here selects
the two-primary subgroup.  Combining (HTS.8)--(HTS.10)
proves that the actual map $M_E\to M_{Q_B}$ is an
isomorphism on those subgroups.

We now retain a known harmonic frame in the lower
full canonical cover.  Use the finite source $D$ and
its full canonical cover $D^c$ from B.2.8 and B.2.26.
The latter result supplies a literal pure element
$b\in D^c$ with
\[
                    [b,s]=1,\qquad
                    \zeta(b)=(12)(34)
                    \quad\text{in the marked }A_4.
 \tag{HTS.11}
\]
The complete polycyclic presentation and the saved
vector in \path{check_source_centralizer_affine.g}
certify these equations in the full cover.  The
source $D$ retains the marked quotient $Q_B$, so its
kernel $K_D$ in $B_0$ is contained in $K$.
There is therefore the actual marked map
\[
 D^c=B_0/[B_0,K_D]\longrightarrow
                   B_0/[B_0,K]=Q_B^c .
 \tag{HTS.12}
\]
Let $\widetilde b_Q$ be the image of $b$ and $b_Q$
its value in $Q_B$.  It is pure, commutes with $s$
already in $Q_B^c$, and its image at projective level
three is the named harmonic involution.  In
particular the canonical commutator obstruction at
$s$ is zero, with no passage to an elementary
central quotient.  The additional replay
\path{check_modular_harmonic_base_frame.g} retains
this exact marked quotient and the saved frame.
It completes the literal dictionary as follows, with
$[a,b]=aba^{-1}b^{-1}$ and $z=(xy)^{-1}$:
\[
 b_Q=q_B(w_4),\qquad w_4=[[x^4,y^4],z^4],\qquad
 \theta_{12}(w_4)=
       \begin{pmatrix}9&4\\8&9\end{pmatrix}.
 \tag{HTS.19}
\]
In particular $b_Q$ has order four; its signed value
is $I$ modulo four and $S$ modulo three.  This is
the result of the actual $D^c\to Q_B$ calculation.
The commuting lift in $Q_B^c$ is the image under
(HTS.12), so the word dictionary alone is not being
used as a central-cover lift.

We lift its value to a centralizer in $E$.  Choose
any pure $u\in E$ over $b_Q$; purity follows from
the common actual quotient $E\to Q_B\to S_3$.
The two involutions $s$ and $usu^{-1}$ lie in the
subgroup $L\langle s\rangle$ of order $2|L|$.
Its Sylow two-subgroups have order two and are
conjugate.  Writing a conjugator as $ns^e$ with
$n\in L$ shows that they are conjugate by an
element of $L$ itself.  Thus, after replacing $u$
by $nu$ for the appropriate $n\in L$, we obtain
\[
               \beta\in P\cap C_E(s),
               \qquad \beta\longmapsto b_Q.
 \tag{HTS.13}
\]
The correction is pure since $L$ lies in $P$.

Choose any pure lift $\widetilde\beta$ of $\beta$
to $E^c$ and put
$d=[\widetilde\beta,s]\in M_E$.
It is central, and $s^2=1$ gives $d^2=1$.
Its image in $M_{Q_B}$ is zero: the image of
$\widetilde\beta$ differs from $\widetilde b_Q$
by a central element, which does not alter the
commutator with $s$.  The injectivity in (HTS.10)
now gives $d=1$.  We have proved the full-cover
equation
\[
                         [\widetilde\beta,s]=1
                         \quad\text{in }E^c .
 \tag{HTS.14}
\]

Apply the exact finite-certificate lifting equality
B.2.7, in its general source form B.2.9, to the
identity automorphism of $E^c$ with pure torsion
frames $(\widetilde\beta,1)$ and cusp frame one.
It supplies an actual full positive automorphism
$\Phi$ of $\widehat B_0$ fixing $\tau=sr$ and
retaining these $E$-values.  Before normalization,
write $f=s\Phi(s)$ and let its harmonic and cubic
frames be $\beta_\Phi,A_\Phi$.  Thus
\[
       f_P=1,\qquad (\beta_\Phi)_P=\beta,
                   \qquad(A_\Phi)_P=1.
 \tag{HTS.15}
\]
If $\beta_\Phi^{\mathrm{ab}}=(u,v)$ and $t=v-u$,
then $f^{\mathrm{ab}}=(t,-t)$, directly from
$f=s\beta_\Phi s\beta_\Phi^{-1}$.
By (HTS.4) and (HTS.15), $t\in18\widehat{\mathbb Z}$.
The exact normalization is
\[
 \Phi_0=\operatorname{Inn}(x^t)\Phi,
 \qquad f_0=y^t f x^{-t},\qquad
 \beta_0=x^t\beta_\Phi,
 \qquad A_0=x^t A_\Phi.
 \tag{HTS.16}
\]
It fixes $\tau$, has $f_0^{\mathrm{ab}}=0$, and
defines a full $g=(1,f_0)\in H$.
Since $x_P,y_P$ have order eighteen, (HTS.16)
preserves all three literal values in (HTS.15).
This proves the finite identities asserted in
(HTS.5), with the actual normalized parameter.

Finally, the signed map $P\to\theta_{12}(F_0)$ is
well-defined because $J_0\subset K$ and $K$ is the
signed kernel on $F_0$.  Therefore
$\theta_{12}(A_g)=I$.  The signed covariance gives
$T_g=V_g=U^{8\rho_g}$ centralizing the named cubic
matrix $R_0$.  At three the three possibilities
$I,U^2,U^4$ contain just one such matrix, namely
$I$.  Hence $\rho_g=0$ modulo three, and $V_g=I$
modulo twelve because it is automatically $I$
modulo four.  Thus $T_g=I$ and
$W_g=\theta_{12}(\beta_g)^{-1}$.
At three the latter has projective value $S$,
by (HTS.11)--(HTS.13); its signed value is therefore
$S$ or $-S$.  Both have harmonic quadratic bit one.
Consequently
\[
       \epsilon_3(T_g)=0,\qquad
       \epsilon_2(W_g)=1,\qquad \Delta(g)=1.
 \tag{HTS.17}
\]
No assertion about the modulo-four component of the
selected harmonic frame is needed.  The five
pentagon insertions occur among the eight maps in
(HTS.3), so every factor of that pentagon is one in
$\Gamma_3$.  This proves the result.
\end{proof}

The fifth retraction of A.4.45 also gives the actual
quotient $P\to F_0/(K'K^3)$.  Let
$J_x=\ker(F_0\to C_2)$ be the parity-of-$x$ subgroup.
The signed modulo-four marking gives $K\subset J_x$,
and therefore $K'K^3\subset J_x'J_x^3$.  Hence
\[
       (f_g)_{F_0/(J_x'J_x^3)}=1,
                  \qquad c_2(f_g)=0\pmod3,
 \tag{HTS.18}
\]
using the actual Schreier coefficient map already
retained in A.4.  This is a further named quotient
of the same genuine source.

The result concerns the ternary quotient (HTS.2).
A.4.45 excludes the earlier identity record with
harmonic frame one and cubic signed frame $5I$;
(HTS.5) uses cubic frame one and the lifted harmonic
frame (HTS.13).  The two conclusions follow from
their displayed actual frame maps.  The additional
binary layer in $N/(N'N^6)$ is not a quotient
retained in this proof.
